\begin{afterword}
\summaryhead

The post answers a common objection to the many-worlds interpretation of quantum mechanics: that its countless unseen worlds violate Occam's razor, which forbids multiplying entities. It calls itself a restatement of ``Belief in the Implied Invisible.''

To calculate simplicity, the author says, there are two roads. One measures a hypothesis by the size of a program that computes its predictions (Kolmogorov complexity and Solomonoff induction); the other by how far it shortens a message describing the data (Minimum Message Length). The post says the two were proven equivalent, and that both count the length of the description, not the memory needed to run it. The conjunction rule explains why each detail a theory must state costs probability. A lottery, a billiards shot and a room's temperature show that objects count against a theory only when they must be specified one by one. History supports this: the universe has kept turning out larger.

A spaceship that passes beyond our cosmological horizon keeps existing, because making it vanish would need a new law. In the same way, many-worlds generates its worlds from compact laws. The post grants that other questions remain, such as why quantum probabilities are not all 50/50, but concludes that the objection from the number of worlds is ``bad math.''

\responsehead

The post's central claim is narrow and, for the formalisms it names, correct. Solomonoff induction and Minimum Message Length charge for the length of a description, not for the number of objects it generates. Its best step is the refinement that objects do count when they must be stated one by one, as in the lottery and billiards cases. This is also the standard reply to the Occam objection: Vaidman's entry on many-worlds in the Stanford Encyclopedia gives it.

The problem is how much the post claims to have settled. The two formalisms are called ``proven equivalent,'' where the author had earlier written ``nearly equivalent'' and named the difference. The promised ``experimental evidence'' is one historical story, for which the author could find no source in the earlier post and neither could I, and a reference to researchers' work that names no study. The formal reason given for ignoring memory and running time does not show that they cannot be counted: Schmidhuber's speed prior counts running time within a normalized prior. Most of all, the last paragraph calls the objection ``entirely defective as probability theory.'' Probability theory does not choose the prior. The post's case for measuring code length rather than objects is strong and widely shared, but it is an argument about which prior to use, not a calculation.

Two further questions bear on simplicity itself. The post names the first, the Born rule, and moves on; it does not raise the second. Vaidman's entry says that getting the right probabilities in many-worlds involves adding a postulate, though it also reports the view that many-worlds is at no disadvantage there. And whether the compact laws alone yield our observations is disputed: Hemmo and Shenker argue that something must be added, and in the formal terms the post uses, Hutter later argued that a program must also locate the observer. None of this shows the objection from the number of worlds to be right. It shows that the post's argument removes one objection and, as the post itself allows, does not decide which theory is simpler overall. That comparison comes later in the book, in ``Many Worlds, One Best Guess.''

Finally, the post uses ``decoherence'' as a name for many-worlds. Physicists use the word for a process every interpretation accepts; the title's claim is about an interpretation.

In short: a correct point about what formal simplicity measures count, presented as more settled than it is: the equivalence of the formalisms is overstated, a dispute about priors is presented as a mathematical error, and the Born rule, which bears on simplicity, is set aside.
\end{afterword}
